\section{Complete Code Examples}

This appendix provides complete, production-ready code implementations for common A/B testing tasks.

\subsection{Comprehensive A/B Test Analysis Class (Python)}

\lstset{language=Python}
\begin{lstlisting}
import numpy as np
import pandas as pd
from scipy import stats
from dataclasses import dataclass
from typing import Optional, Dict, List

@dataclass
class ABTestResult:
    """Container for A/B test results."""
    control_mean: float
    treatment_mean: float
    difference: float
    relative_lift: float
    standard_error: float
    confidence_interval: tuple
    p_value: float
    statistically_significant: bool
    test_statistic: float
    sample_size_control: int
    sample_size_treatment: int

class ABTestAnalyzer:
    """
    Comprehensive A/B test analysis.

    Supports:
    - Proportions (conversion rates)
    - Continuous metrics (revenue, time, etc.)
    - Multiple testing corrections
    - Subgroup analysis
    """

    def __init__(self, alpha=0.05):
        self.alpha = alpha

    def test_proportions(self, control_conversions, control_total,
                        treatment_conversions, treatment_total,
                        method='z_test'):
        """
        Test difference in proportions.

        Parameters:
        -----------
        control_conversions : int
            Number of successes in control
        control_total : int
            Total observations in control
        treatment_conversions : int
            Number of successes in treatment
        treatment_total : int
            Total observations in treatment
        method : str
            'z_test', 'chi_squared', or 'fisher_exact'

        Returns:
        --------
        ABTestResult : Test results
        """
        p1 = control_conversions / control_total
        p2 = treatment_conversions / treatment_total

        if method == 'z_test':
            # Pooled proportion
            p_pool = ((control_conversions + treatment_conversions) /
                     (control_total + treatment_total))

            # Standard error under null
            se_null = np.sqrt(p_pool * (1 - p_pool) *
                            (1/control_total + 1/treatment_total))

            # Test statistic
            z_stat = (p2 - p1) / se_null
            p_value = 2 * (1 - stats.norm.cdf(abs(z_stat)))

            # Confidence interval (unpooled)
            se_unpooled = np.sqrt(p1*(1-p1)/control_total +
                                 p2*(1-p2)/treatment_total)
            z_crit = stats.norm.ppf(1 - self.alpha/2)
            ci = (p2 - p1 - z_crit * se_unpooled,
                  p2 - p1 + z_crit * se_unpooled)

            test_stat = z_stat
            se = se_unpooled

        elif method == 'chi_squared':
            table = [[control_conversions, control_total - control_conversions],
                    [treatment_conversions, treatment_total - treatment_conversions]]
            chi2, p_value, dof, expected = stats.chi2_contingency(table)
            test_stat = chi2

            # CI from z-test
            se = np.sqrt(p1*(1-p1)/control_total +
                        p2*(1-p2)/treatment_total)
            z_crit = stats.norm.ppf(1 - self.alpha/2)
            ci = (p2 - p1 - z_crit * se, p2 - p1 + z_crit * se)

        elif method == 'fisher_exact':
            table = [[control_conversions, control_total - control_conversions],
                    [treatment_conversions, treatment_total - treatment_conversions]]
            _, p_value = stats.fisher_exact(table)
            test_stat = None

            # CI from z-test
            se = np.sqrt(p1*(1-p1)/control_total +
                        p2*(1-p2)/treatment_total)
            z_crit = stats.norm.ppf(1 - self.alpha/2)
            ci = (p2 - p1 - z_crit * se, p2 - p1 + z_crit * se)

        return ABTestResult(
            control_mean=p1,
            treatment_mean=p2,
            difference=p2 - p1,
            relative_lift=(p2 - p1) / p1 if p1 > 0 else np.nan,
            standard_error=se,
            confidence_interval=ci,
            p_value=p_value,
            statistically_significant=p_value < self.alpha,
            test_statistic=test_stat,
            sample_size_control=control_total,
            sample_size_treatment=treatment_total
        )

    def test_means(self, control_data, treatment_data,
                  equal_variance=False):
        """
        Test difference in means.

        Parameters:
        -----------
        control_data : array-like
            Control group values
        treatment_data : array-like
            Treatment group values
        equal_variance : bool
            Assume equal variances (default False, uses Welch's t-test)

        Returns:
        --------
        ABTestResult : Test results
        """
        control_data = np.array(control_data)
        treatment_data = np.array(treatment_data)

        n1 = len(control_data)
        n2 = len(treatment_data)
        mean1 = np.mean(control_data)
        mean2 = np.mean(treatment_data)
        var1 = np.var(control_data, ddof=1)
        var2 = np.var(treatment_data, ddof=1)

        # T-test
        t_stat, p_value = stats.ttest_ind(treatment_data, control_data,
                                          equal_var=equal_variance)

        # Standard error and CI
        se = np.sqrt(var1/n1 + var2/n2)

        if equal_variance:
            df = n1 + n2 - 2
        else:
            # Welch-Satterthwaite degrees of freedom
            df = ((var1/n1 + var2/n2)**2 /
                 ((var1/n1)**2/(n1-1) + (var2/n2)**2/(n2-1)))

        t_crit = stats.t.ppf(1 - self.alpha/2, df)
        ci = (mean2 - mean1 - t_crit * se, mean2 - mean1 + t_crit * se)

        return ABTestResult(
            control_mean=mean1,
            treatment_mean=mean2,
            difference=mean2 - mean1,
            relative_lift=(mean2 - mean1) / mean1 if mean1 > 0 else np.nan,
            standard_error=se,
            confidence_interval=ci,
            p_value=p_value,
            statistically_significant=p_value < self.alpha,
            test_statistic=t_stat,
            sample_size_control=n1,
            sample_size_treatment=n2
        )

    def calculate_sample_size_proportions(self, p1, p2, power=0.80):
        """
        Calculate required sample size for proportion test.

        Parameters:
        -----------
        p1 : float
            Control proportion
        p2 : float
            Treatment proportion
        power : float
            Desired statistical power

        Returns:
        --------
        dict : Sample size requirements
        """
        z_alpha = stats.norm.ppf(1 - self.alpha/2)
        z_beta = stats.norm.ppf(power)

        variance = p1*(1-p1) + p2*(1-p2)
        effect = (p2 - p1)**2

        n_per_group = int(np.ceil(
            ((z_alpha + z_beta)**2 * variance) / effect
        ))

        return {
            'n_per_group': n_per_group,
            'n_total': 2 * n_per_group,
            'p1': p1,
            'p2': p2,
            'effect_absolute': p2 - p1,
            'effect_relative': (p2 - p1) / p1,
            'alpha': self.alpha,
            'power': power
        }

    def format_results(self, result: ABTestResult) -> str:
        """Format results for reporting."""
        report = f"""
A/B Test Results
================
Control Mean: {result.control_mean:.4f}
Treatment Mean: {result.treatment_mean:.4f}

Difference: {result.difference:.4f}
Relative Lift: {result.relative_lift:.2%}
Standard Error: {result.standard_error:.4f}

95% Confidence Interval: [{result.confidence_interval[0]:.4f}, {result.confidence_interval[1]:.4f}]

Test Statistic: {result.test_statistic:.4f}
P-value: {result.p_value:.4f}
Statistically Significant: {'Yes' if result.statistically_significant else 'No'}

Sample Sizes:
  Control: {result.sample_size_control:,}
  Treatment: {result.sample_size_treatment:,}
"""
        return report

# Example usage
if __name__ == "__main__":
    analyzer = ABTestAnalyzer(alpha=0.05)

    # Test proportions
    result_prop = analyzer.test_proportions(
        control_conversions=80,
        control_total=1000,
        treatment_conversions=95,
        treatment_total=1000
    )

    print(analyzer.format_results(result_prop))

    # Calculate sample size
    sample_size = analyzer.calculate_sample_size_proportions(
        p1=0.08,
        p2=0.09,
        power=0.80
    )

    print("\nSample Size Calculation:")
    print(f"  Required per group: {sample_size['n_per_group']:,}")
    print(f"  Total required: {sample_size['n_total']:,}")
\end{lstlisting}

\subsection{Experiment Analysis Pipeline (R)}

\lstset{language=R}
\begin{lstlisting}
library(tidyverse)
library(broom)

# Comprehensive experiment analysis pipeline
analyze_experiment <- function(data, metric, variant_col = "variant",
                              alpha = 0.05) {
  #' Analyze A/B test experiment
  #'
  #' @param data DataFrame with experiment data
  #' @param metric Character, name of metric column
  #' @param variant_col Character, name of variant column
  #' @param alpha Numeric, significance level
  #'
  #' @return List with analysis results

  # 1. Pre-analysis checks
  cat("=== Pre-Analysis Checks ===\n")

  # Sample ratio mismatch
  variant_counts <- table(data[[variant_col]])
  srm_test <- chisq.test(variant_counts)
  cat(sprintf("SRM p-value: %.4f %s\n", srm_test$p.value,
              ifelse(srm_test$p.value < 0.01, "(WARNING)", "(OK)")))

  # 2. Summary statistics
  cat("\n=== Summary Statistics ===\n")
  summary_stats <- data %>%
    group_by(!!sym(variant_col)) %>%
    summarise(
      n = n(),
      mean = mean(!!sym(metric), na.rm = TRUE),
      median = median(!!sym(metric), na.rm = TRUE),
      sd = sd(!!sym(metric), na.rm = TRUE),
      se = sd / sqrt(n),
      .groups = 'drop'
    )

  print(summary_stats)

  # 3. Statistical test
  cat("\n=== Statistical Test ===\n")

  variants <- unique(data[[variant_col]])
  control <- data %>% filter(!!sym(variant_col) == variants[1]) %>%
    pull(!!sym(metric))
  treatment <- data %>% filter(!!sym(variant_col) == variants[2]) %>%
    pull(!!sym(metric))

  # Welch's t-test (robust to unequal variances)
  test_result <- t.test(treatment, control, var.equal = FALSE)

  cat(sprintf("Test statistic: %.4f\n", test_result$statistic))
  cat(sprintf("P-value: %.4f\n", test_result$p.value))
  cat(sprintf("Significant: %s\n",
              ifelse(test_result$p.value < alpha, "Yes", "No")))

  # 4. Effect sizes
  cat("\n=== Effect Sizes ===\n")

  mean_control <- mean(control, na.rm = TRUE)
  mean_treatment <- mean(treatment, na.rm = TRUE)
  diff <- mean_treatment - mean_control
  relative_lift <- diff / mean_control

  cat(sprintf("Absolute difference: %.4f\n", diff))
  cat(sprintf("Relative lift: %.2f%%\n", relative_lift * 100))
  cat(sprintf("95%% CI: [%.4f, %.4f]\n",
              test_result$conf.int[1], test_result$conf.int[2]))

  # 5. Visualizations
  p1 <- ggplot(data, aes_string(x = variant_col, y = metric)) +
    geom_boxplot() +
    geom_jitter(alpha = 0.1, width = 0.2) +
    labs(title = "Distribution by Variant",
         x = "Variant", y = metric) +
    theme_minimal()

  # Density plot
  p2 <- ggplot(data, aes_string(x = metric, fill = variant_col)) +
    geom_density(alpha = 0.5) +
    labs(title = "Density Plot by Variant",
         x = metric, fill = "Variant") +
    theme_minimal()

  # Return results
  list(
    srm_check = list(p_value = srm_test$p.value,
                    passed = srm_test$p.value >= 0.01),
    summary = summary_stats,
    test = test_result,
    effect_size = list(
      absolute = diff,
      relative = relative_lift,
      ci = test_result$conf.int
    ),
    plots = list(boxplot = p1, density = p2)
  )
}

# Bayesian analysis function
bayesian_ab_test <- function(x_control, n_control, x_treatment, n_treatment,
                             prior_alpha = 1, prior_beta = 1,
                             n_sim = 100000) {
  #' Bayesian A/B test for proportions
  #'
  #' @param x_control Conversions in control
  #' @param n_control Sample size control
  #' @param x_treatment Conversions in treatment
  #' @param n_treatment Sample size treatment
  #' @param prior_alpha Beta prior alpha
  #' @param prior_beta Beta prior beta
  #' @param n_sim Number of Monte Carlo simulations
  #'
  #' @return List with Bayesian results

  # Posterior parameters
  post_alpha_control <- prior_alpha + x_control
  post_beta_control <- prior_beta + n_control - x_control
  post_alpha_treatment <- prior_alpha + x_treatment
  post_beta_treatment <- prior_beta + n_treatment - x_treatment

  # Simulate from posteriors
  p_control <- rbeta(n_sim, post_alpha_control, post_beta_control)
  p_treatment <- rbeta(n_sim, post_alpha_treatment, post_beta_treatment)

  # Probability treatment > control
  prob_treatment_better <- mean(p_treatment > p_control)

  # Expected lift
  lift <- p_treatment / p_control - 1
  expected_lift <- mean(lift)
  lift_ci <- quantile(lift, c(0.025, 0.975))

  # Expected loss
  loss_control <- mean(pmax(0, p_treatment - p_control))
  loss_treatment <- mean(pmax(0, p_control - p_treatment))

  list(
    prob_treatment_better = prob_treatment_better,
    expected_lift = expected_lift,
    lift_ci = lift_ci,
    loss = list(
      control = loss_control,
      treatment = loss_treatment,
      recommendation = ifelse(loss_control < loss_treatment,
                             "Control", "Treatment")
    )
  )
}

# Example usage:
# data <- read_csv("experiment_data.csv")
# results <- analyze_experiment(data, metric = "conversion")
# print(results$plots$boxplot)
\end{lstlisting}

\subsection{Experiment Simulation Framework}

\lstset{language=Python}
\begin{lstlisting}
import numpy as np
import pandas as pd
from typing import Callable

class ExperimentSimulator:
    """Simulate A/B tests for power analysis and methodology validation."""

    def __init__(self, seed=None):
        if seed is not None:
            np.random.seed(seed)

    def simulate_proportions(self, n_control, n_treatment,
                            p_control, p_treatment,
                            n_simulations=1000, alpha=0.05):
        """
        Simulate A/B tests with binary outcomes.

        Returns:
        --------
        dict : Simulation results including empirical power and Type I error
        """
        results = []

        for _ in range(n_simulations):
            # Generate data
            control = np.random.binomial(1, p_control, n_control)
            treatment = np.random.binomial(1, p_treatment, n_treatment)

            # Test
            p1 = np.mean(control)
            p2 = np.mean(treatment)

            p_pooled = (np.sum(control) + np.sum(treatment)) / (n_control + n_treatment)
            se = np.sqrt(p_pooled * (1 - p_pooled) * (1/n_control + 1/n_treatment))

            z = (p2 - p1) / se
            p_value = 2 * (1 - stats.norm.cdf(abs(z)))

            results.append({
                'p_value': p_value,
                'reject': p_value < alpha,
                'p1_hat': p1,
                'p2_hat': p2,
                'effect_estimate': p2 - p1
            })

        results_df = pd.DataFrame(results)

        # Calculate metrics
        empirical_power = results_df['reject'].mean()
        type_i_error = empirical_power if p_control == p_treatment else None

        return {
            'results': results_df,
            'empirical_power': empirical_power,
            'type_i_error': type_i_error,
            'mean_p_value': results_df['p_value'].mean(),
            'significant_count': results_df['reject'].sum()
        }

# Example: Verify power calculation
simulator = ExperimentSimulator(seed=42)
sim_result = simulator.simulate_proportions(
    n_control=1000,
    n_treatment=1000,
    p_control=0.10,
    p_treatment=0.12,
    n_simulations=10000,
    alpha=0.05
)

print(f"Empirical power: {sim_result['empirical_power']:.3f}")
print(f"Significant tests: {sim_result['significant_count']}/{10000}")
\end{lstlisting}
